%% Use the "normalphoto" option if you want a normal photo instead of cropped to a circle
% \documentclass[10pt,a4paper,withhyper,normalphoto]{altacv}

\documentclass[8pt, a4paper, withhyper]{altacv}
%% AltaCV uses the fontawesome5 and simpleicons packages.
%% See http://texdoc.net/pkg/fontawesome5 and http://texdoc.net/pkg/simpleicons for full list of symbols.

% Change the page layout if you need to
\geometry{left=1.25cm, right=1.25cm, top=1.5cm, bottom=1.5cm, columnsep=1.2cm}

% The paracol package lets you typeset columns of text in parallel
\usepackage{paracol}

\iftutex
% If using xelatex or lualatex:
\setmainfont{Roboto Slab}
\setsansfont{Lato}
\renewcommand{\familydefault}{\sfdefault}
\else
% If using pdflatex:
\usepackage[rm]{roboto}
\usepackage[defaultsans]{lato}
% \usepackage{sourcesanspro}
\renewcommand{\familydefault}{\sfdefault}
\fi

% GitHub-like dark theme (refined)
\definecolor{GitHubBg}{HTML}{0D1117}        % main background
\definecolor{GitHubSurface}{HTML}{161B22}   % slightly lighter panels
\definecolor{GitHubBorder}{HTML}{30363D}    % borders / separators
\definecolor{GitHubText}{HTML}{C9D1D9}      % primary text
\definecolor{GitHubMuted}{HTML}{8B949E}     % secondary text
\definecolor{GitHubAccent}{HTML}{58A6FF}    % links / highlights (blue)
\definecolor{GitHubGreen}{HTML}{3FB950}     % subtle green (not neon)
\definecolor{GitHubRed}{HTML}{F85149}       % optional accent

% Apply colors
\colorlet{name}{GitHubText}
\colorlet{tagline}{GitHubAccent}
\colorlet{heading}{GitHubAccent}
\colorlet{headingrule}{GitHubBorder}
\colorlet{subheading}{GitHubText}
\colorlet{accent}{GitHubGreen}
\colorlet{emphasis}{GitHubText}
\colorlet{body}{GitHubText}

\pagecolor{GitHubBg}
\color{GitHubText}

% Change some fonts, if necessary
\renewcommand{\namefont}{\Huge\rmfamily\bfseries}
\renewcommand{\personalinfofont}{\footnotesize}
\renewcommand{\cvsectionfont}{\LARGE\rmfamily\bfseries}
\renewcommand{\cvsubsectionfont}{\large\bfseries}

% Change the bullets for itemize and rating marker
% for \cvskill if you want to
\renewcommand{\cvItemMarker}{{\small\textbullet}}
\renewcommand{\cvRatingMarker}{\faCircle}
% ...and the markers for the date/location for \cvevent
% \renewcommand{\cvDateMarker}{\faCalendar*[regular]}
% \renewcommand{\cvLocationMarker}{\faMapMarker*}

\input{pubs-num}

%% sample.bib contains your publications
\addbibresource{sample.bib}

\begin{document}
  \name{Matthew de Lannoy} \tagline{Mechanical Engineer}
  \photoR{2.8cm}{profile}

  \personalinfo{%

  % Not all of these are required!
  \email{matthewdelan@gmail.com} \phone{+31 6 45035492} \location{Delft, The Netherlands}
  % \homepage{www.homepage.com}

  % \linkedin{your_id}

  \github{machine0herald}
  }

  \makecvheader
  \AtBeginEnvironment{itemize}{\small}

  %% Set the left/right column width ratio to 6:4.
  \columnratio{0.6}

  % Start a 2-column paracol. Both the left and right columns will automatically
  % break across pages if things get too long.
  \begin{paracol}
    {2}
    % \cvsection{Experience}

    % \cvevent{Job Title 1}{Company 1}{Month 20XX -- Ongoing}{Location}
    % \begin{itemize}
    % \item Job description 1
    % \item Job description 2
    % \end{itemize}

    % \divider

    % \cvevent{Job Title 2}{Company 2}{Month 20XX -- Ongoing}{Location}
    % \begin{itemize}
    % \item Job description 1
    % \item Job description 2
    % \end{itemize}

    % ////////////////////////////////////////////////////////////////////

    \cvsection{Projects}
    
    \cvevent{
      \href{https://github.com/C-O-R-A}{Cora} 
      \href{https://github.com/C-O-R-A}{\faGithub}
      \href{https://cad.onshape.com/documents/971c71e558f4d758be7c1282/w/a52d355fced062b7019a3a28/e/e9a08b2016a881ad8b69894a}{\faCube\ CAD}
    }{Robotics Student Association}{Mar 2025 -- Present} {}    
    \textbf{Cora} is a 6 degree-of-freedom cobot (Collaborative Robot arm) that I am
    designing and building as a member of the Robotics Student Association (RSA)
    at TU Delft. In this project I have taken on the role of Lead Mechanical,
    Electrical and Software Engineer. Since March 2026, the team has expanded 
    to five members with backgrounds in Mechanical Engineering, 
    Electrical Engineering and Computer Science.
    The cobot is intended as a platform for educational and research purposes.
    The robot has a reach of 1.2m and a maximum payload of 4kg.
    It runs ROS2 Jazzy on an onboard Raspberry Pi 5 SBC
    and also features a socket interface and Python SDK for running programs 
    remotely through LAN or WiFi.
    
    \divider
    
    \cvevent{
      \href{https://github.com/C-O-R-A/cora-chess}{You vs Cobot} 
      \href{https://github.com/C-O-R-A/cora-chess}{\faGithub}
      }{Robotics Student Association}{Mar 2025 -- Present} {}
      Running in parallel with \textbf{Cora} is its first application, \textbf{You vs Cobot}, 
      where the Ethernet interface and Python SDK are used 
      to make the robot play chess against a human opponent.
      
      \divider
      
      \cvevent{Cavolab}{Prysmian / Robohouse}{Sept 2025 - Present}{} As part of my
      minor, I worked on a project to design and build an automated cable cleaning
      robot for high voltage cable manufacturer Prysmian. 
      The project was carried out in collaboration with
      Robohouse, who provide expertise in automation and robotics as well as a workshop
      and office space. My role in the team was that of Lead Mechanical Engineer. 
      As part of that role, I was responsible for designing the linear gantry and
      cable clamps as well as the structures that secure several pistons for securing and
      working on the cable.
      
      \divider
      
      \cvevent{Atlas}{Robotics Student Association}{Sept 2024 - Mar 2025}{} Atlas
      is a mostly 3D printed 6-DOF desktop robotic arm that I designed as part of
      the Robotics Student Association. It has a reach of 500mm and a maximum payload
      of 1kg. The arm was able to be controlled through the onboard computer. I later
      decided to create a new arm under the name 'Cora', keeping in mind
      the new insights I had gained.
      
      \divider
      
      \cvevent{Hexapod}{Robotics Student Association}{May 2024}{} In May 2024,
      I designed and built a hexapod robot. This robot was designed to be used as
      a platform for testing various sensors and algorithms for autonomous
      navigation. The robot features six legs, each with three degrees of freedom,
      and is capable of walking and climbing over obstacles. The robot is
      controlled using a Servo2040 which has an RP2040 chip and is programmed
      using Python.
      
      % ////////////////////////////////////////////////////////////////////
      
      
      \cvsection{Education} 
      \cvevent{B.Sc.\ Mechanical Engineering}{TU Delft}{Sept 2022 -- June 2026}{}
      \cvevent{Minor\ Robotics}{TU Delft}{Sept 2025 -- Jan 2026}{}
      
      \switchcolumn
      
      % ////////////////////////////////////////////////////////////////////

      \cvsection{Strengths and Skills}
      \smallskip
      
      \cvsubsection{Programming} \cvtag{C++} \cvtag{Python} \cvtag{ROS 2} \cvtag{Platform IO}
      \cvtag{Git}\\
      
      \cvsubsection{Manufacturing} \cvtag{3D Printing} \cvtag{Sheetmetal Part Design}
      \cvtag{Design for CNC machining} \cvtag{PCB Design}\\
      
      \cvsubsection{CAD} \cvtag{Solidworks} \cvtag{Onshape} \cvtag{KiCAD}\\
      % \cvtag{Embedded Systems}
      
    
      % ////////////////////////////////////////////////////////////////////
    \cvsection{Extracurricular} \cvevent{Workshop Committee}{Robotics Student Association Delft}{Sept 2025 -- Current}{}
    \cvevent{Student Mentor}{TU Delft - Faculty of Mechanical Engineering}{Sept 2024 -- Feb 2025}{}
    \cvevent{Teaching Assistant - CAD}{TU Delft - Faculty of Mechanical Engineering}{Nov 2023 -- Feb 2026}{}
    \cvevent{Teaching Assistant - Robotics}{TU Delft - Faculty of Mechanical Engineering}{Feb 2026 -- Apr 2026}{}
    
    \medskip
    
    % \cvsection{References}
    
    % \cvref{name}{email}{mailing address}
    % \cvref{Prof.\ Werner van de Sande}{TU Delft - PME}{a.beta@university.edu}
    
    % \divider
    
    % \cvref{Prof.\ Martijn Wisse}{TU Delft - COR}{TU Delft - COR}
    
    % \divider
    
    % \cvref{Prof.\ Martin Klomp}{Robohouse}{g.delta@university.edu}
    
    % ////////////////////////////////////////////////////////////////////


    \cvsection{Certifications}
    \begin{small}
      \cvsubsection{SolidWorks} \cvevent{CSWA (Associate)}{}{}{}
        \cvevent{CSWA – Simulation}{}{}{}
        \cvevent{CSWP (Professional)}{}{}{}
        \cvevent{CSWPA (Professional Advanced) – Drawings}{}{}{}
        \cvevent{CSWPA – Weldments}{}{}{}
        \cvevent{CSWPA – Sheetmetal}{}{}{}
        \cvevent{CSWPA – Surfacing}{}{}{}
        \cvevent{CSWE (Expert)}{}{}{}
    \end{small}

  \end{paracol}
  \end{document}